Formal methods have been increasingly applied to verify the correctness of compiler optimizations,
ensuring that transformations preserve program semantics.
This report describes the implementation of \proj,
a \lean formalization of a small executable \llvm-like intermediate
representation (\ir) with symbolic integer widths, a surface syntax embedded in
\lean, operational semantics, structural well-formedness checks, and
machine-checked refinement proofs.

Two examples illustrate the result. First, the project proves that
\texttt{add x, 0} is refined by \texttt{return x} both at a fixed width and for all
instantiations of a symbolic width. Second, the project models \vundef using
an explicit supply of values and proves a width-generic theorem relating
\texttt{add undef, undef} and \texttt{mul undef, 2}, showing that the latter
is a more defined target behavior than the former. These examples capture
what the project ultimately achieved well: not a broad model of all of
\llvm, but a small, readable, executable framework in which subtle semantic
issues such as symbolic widths, \poison, and \vundef can already be studied
with machine-checked proofs.
